% LSPV-3::BRAUN-TO-MIXED-NONPROJECTION369::2026-11-03_0924_PST
\documentclass[11pt]{article}
\usepackage[margin=0.88in]{geometry}
\usepackage{amsmath,amssymb,amsthm,hyperref}
\newtheorem{theorem}{Theorem}
\newtheorem{lemma}{Lemma}
\title{LSPV-3 369: Quantitative Impossibility of Backtranslating Short Mixed BPHP Certificates to Resolution over Parities}
\author{PNP Hypergraph \& Proof Certificates, LSPV-3}
\date{November 3, 2026, 9:24 AM PST}
\begin{document}\maketitle
\noindent\textbf{Proof status.} \textbf{PROVED/EXACT CONDITIONAL on the external Braun (2026) theorem} for the new quantitative nonprojection; the short mixed certificate bound was already \textbf{PROVED/EXACT} and computationally checked by this same project in LSPV-3 346, not rediscovered. \textbf{COMPUTATIONAL} for finite signed source and parity controls. Braun is a new, AI-assisted, Lean-formalized \emph{preprint claim}; this run did not independently replay its Lean development. No unconditional superpolynomial lower bound for the full mixed calculus is asserted. B1, B2, B3 and P versus NP remain \textbf{OPEN}.

\section*{Fixed system, quantifiers, and published starting point}
Keep the exact preselected source-authenticated, binary-premise, unit-weight $R_{\mathrm{mix341}}$: RES, AND3, GF$(2)$ row addition, justified LOCAL-CONVERT and one charged conservative signed two-input VARDEF. A VARDEF produces three authenticated clauses, is acyclic, and cannot be an instance-specific answer axiom. Every derived inference node has unit weight, source clauses cost zero, sharing counts once, and optimal adaptive legal proof selection is free. For valid general encoded CNF $F$ and length $N$ in \emph{bits},
\[
 W_R(F)=\min_{\pi:F\vdash_R\bot}|\pi_{\rm derived}|,\qquad
 \overline W_R(N)=\max_{\substack{F\ {\rm valid\ UNSAT}\\|\operatorname{enc}(F)|\le N}}W_R(F).
\]
The empty maximum is zero. Ordinary RES makes the set of attainable nonnegative integer costs nonempty for UNSAT $F$, so its minimum is attained. This is a proof-DAG cost, not the running time of a SAT algorithm or a checker.

For $\ell\ge 2$ let $n=2^\ell$ and let $F_\ell$ be the \emph{original} bit pigeonhole CNF: pigeon $i\in[n+1]$ has $\ell$ bits $x_{i,b}$ and, for every pair $i<j$ and hole word $h\in\{0,1\}^{\ell}$, there is a collision clause
\[
 D_{ijh}=\bigvee_{b=1}^\ell (x_{i,b}\ne h_b)\ \vee\
          \bigvee_{b=1}^\ell (x_{j,b}\ne h_b).
\]
It has $(n+1)\ell$ original variables, $m_\ell=n\binom{n+1}{2}$ indispensable width-$2\ell$ sources, and $N_\ell=\Theta(n^3\ell^2)$ explicit signed ASCII DIMACS bits (not merely $\Theta(n^3\ell)$ literal occurrences). Assigning the same hole $h$ to $i,j$ and bijecting the others to the remaining holes gives the SAT witness for deletion or any one literal sign reversal of $D_{ijh}$. All $m_\ell$ sources are needed in a sound conservative proof, so the binary-premise ancestry bound gives $W_R(F_\ell)\ge m_\ell-1$.

\noindent\textbf{External lower bound, MODEL-RESTRICTED.} Braun's September 2026 preprint [1] claims, with an accompanying Lean 4 formalization, for every integer $\ell\ge32$, every unrestricted DAG-like Resolution-over-Parities (Res$(\oplus)$) refutation $\rho$ of the \emph{same} $F_\ell$, without regularity or depth restrictions, has a number $S_\oplus(\rho)$ of all line nodes, \emph{including original axioms}, satisfying
\[
 S_\oplus(\rho)>\exp\!\left(\frac{n}{32768\ell^2}\right).
 \tag{1}
\]
Res$(\oplus)$ clauses are disjunctions of affine GF$(2)$ equations. The theorem is about this fixed proof model, not general conservative extensions, arbitrary circuit reasoning, machine time, or the full $R_{\mathrm{mix341}}$. We use equation (1) as a clearly identified external theorem hypothesis, not as something independently verified by the finite controls.

\noindent\textbf{Authenticated project upper bound, PRIOR result.} LSPV-3 BPHP-FULL-COOK-CERTIFICATE346 (November 2, 2026, 16:24 PST; Git blob \texttt{e4931b197b759708fc53983c028894dc35c1eaab}) gives, for every $\ell\ge2$, a uniform authentic source-preserving RES/DEF refutation $\pi_\ell$ of $F_\ell$, using a prefix-indicator translation to unary PHP, followed by a charged local Cook elimination. Its derived-node count $q_\ell$ satisfies
\[
 m_\ell-1\le q_\ell\le C(n^4+n^3\ell^2)\le C_0 n^4
 \quad(n=2^\ell\ge4)
 \tag{2}
\]
for fixed constants $C,C_0>0$. Its complete $n=4,8,16$ proof DAGs were previously replayed independently, including every signed pivot and gate slot. The earlier source listed respective charged totals 590, 8,073 and 112,795; these are \emph{upper} certificates, not minima. The present work neither repeats that compiler nor treats its discovery cost as zero.

\section*{New quantitative transfer obstruction}
\begin{theorem}[Same-source backtranslation output lower bound]
Assume Braun's inequality (1) applies with the stated original bit-PHP clauses and complete unrestricted DAG-like Res$(\oplus)$ rules. For every $\ell\ge32$, choose the pre-existing certified mixed proof $\pi_\ell$ satisfying (2). For \emph{every} valid Res$(\oplus)$ refutation $\rho_\ell$ of that same $F_\ell$, regardless of whether $\rho_\ell$ is returned by a uniform algorithm, nonuniform family, computationally unbounded map or an omniscient after-the-fact selector,
\[
 S_\oplus(\rho_\ell)>
 \exp\!\left(\frac{9}{32768 C_0^{1/4}}\,
 \frac{q_\ell^{1/4}}{(\log_2(q_\ell+1)+1)^2}\right).
 \tag{3}
\]
Consequently \emph{no} same-source polynomial-output backtranslation exists even on just this uniformly explicit family of polynomial-cost $R_{\mathrm{mix341}}$ proofs. Charging only derived Res$(\oplus)$ nodes changes the lower bound by subtracting $m_\ell$, negligible asymptotically. This is a \emph{family-local output-size obstruction}, not a SAT-runtime theorem.
\end{theorem}
\begin{proof}
By (2), $n\ge(q_\ell/C_0)^{1/4}$. Since all $m_\ell$ source clauses are indispensable,
$q_\ell+1\ge m_\ell=n^2(n+1)/2\ge n^3/2$, so
$\ell=\log_2 n\le(\log_2(q_\ell+1)+1)/3$.
Inserting these two inequalities into the right-hand exponent of (1) proves (3). Its exponent eventually dominates $\log(q_\ell+2)$, ruling out every fixed polynomial output size in $q_\ell+N_\ell$ since $N_\ell=\Theta(n^3\ell^2)$ and $q_\ell\ge\Theta(n^3)$. The original $m_\ell$ source lines are only polynomially many, so subtracting them does not affect the superpolynomial conclusion. The proof makes \emph{no} statement that constructing $\rho_\ell$ costs exactly $S_\oplus$, nor any claim about the fastest correct SAT decider.
\end{proof}

\noindent\textbf{Direction firewall.} $R_{\mathrm{mix341}}$ \emph{contains} charged ER2, so $W_{R_{\mathrm{mix341}}}(F_\ell)\le q_\ell=\operatorname{poly}(N_\ell)$. A hard restricted Res$(\oplus)$ proof can be exponentially longer on this \emph{identical} source. Thus its hard lower bound cannot be copied to the minimum over the larger fixed system; the smaller-to-larger simulation inequality goes the opposite way. This is not the first project proof of nonprojection: BPHP343 and BPHP346, LSPV-2 HAKEN-ER2-NONPROJECTION364, and LSPV-5 REFEREE316 already ruled out polynomial-cost backward transfers for their respective hard families. Equation (3) refines the Res$(\oplus)$ output-obstruction rate explicitly \emph{as a function of the short mixed proof node count}, with a publication-status qualification on the new preprint.

\section*{Target bridges and four precise falsifiers}
A fixed total correct SAT decider $A$ handles SAT, UNSAT and malformed encodings and has
$T_A^{wc}(N)=\max_{x\in\{0,1\}^{\le N}}T_A(x)$.
B1 states $\forall A\,\exists p_A\,\forall F\in\mathrm{UNSAT}$,
$W_R(F)\le p_A(|\operatorname{enc}F|+T_A(F))$.
B2 states $\forall A\,\exists p_A\,\forall N$,
$\overline W_R(N)\le p_A(N+T_A^{wc}(N))$.
B3 states $\mathrm{SAT}\in\mathrm P\Longrightarrow\overline W_R(N)\le\operatorname{poly}(N)$ for this \emph{same} $R$.
B1 implies B2 implies B3. Their proofs are separate and \emph{OPEN}.
\begin{enumerate}
\item[] F1: Every short mixed $F_\ell$ refutation can be translated into a same-source Res$(\oplus)$ refutation with polynomially many lines. \textbf{FALSIFIED/RETIRED} conditional on Braun and by (3), even without any runtime condition on the translator.
\item[] F2: A published unrestricted-DAG Res$(\oplus)$ lower bound forces a superpolynomial $W_{R_{\mathrm{mix341}}}$ on $F_\ell$. \textbf{FALSIFIED/RETIRED} unconditionally by existing authenticated polynomial mixed certificates.
\item[] F3: Equation (3) proves every uniform correct SAT algorithm needs superpolynomial time. \textbf{UNSUPPORTED}: a restricted proof backtranslation output blowup is not a machine lower bound; no B1/B2/B3 follows.
\item[] F4: $T_A(M)\le M^{K+o(1)}$ implies each call to $A$ costs $\Omega(M^{K+o(1)})$. \textbf{FALSIFIED}: an upper bound does not provide a matching lower bound; a solver can return immediately on a recognizable input.
\end{enumerate}

\section*{Independent source and control ledger; costs}
The separate no-import JavaScript checker \texttt{experiments/arc/lspv3_369_braun_backtranslation_controls.js} reconstructed the signed sources for $\ell=1,2,3$, checking counts $(m,v)=(6,3),(40,10),(288,27)$, every deletion and every literal sign flip for the first two sizes, and a deterministic 32-source sample at $\ell=3$. Total checked: 78 deletion SAT witnesses and 364 occurrence sign-flip SAT witnesses, plus matched injective $n$-into-$n$ SAT controls. Full truth-table enumeration for $\ell=1,2$ found zero satisfying assignments. Separate $K_4$ Tseitin parity controls had 8 even-charge solutions and zero odd-charge solutions, expressly preventing an illicit transfer from hard Resolution to inexpensive global XOR. \emph{These finite checks do not verify Braun (1), certify a shortest mixed proof, or rerun the large BPHP346 proofs.} Raw counts are published separately.

Cost separation is strict: source clauses $m_\ell$ and their bytes, derived nodes $q_\ell$, encoded proof evidence (already bounded polynomially in prior authenticated BPHP346), verification work, normalization, actual proof-discovery time and all-phase peak space are different quantities. Equation (3) is a lower bound on \emph{encoded output node count under restricted proof-type conversion}, not on general computation. Nothing infers an FLMV simultaneous time--space contradiction, and the corrected space envelope $s'\le\max\{\rho,\eta,\lambda,\lambda s\}$ continues to charge every phase via $\rho$.

\section*{Repository, literature, assessments and write safeguards}
The November 2, 12:04 PST five-role directive was read first from the requested current GitHub path (user-pinned historical commit \texttt{4e29ed7}, blob \texttt{033b54b...}); then current issue \#3, root \texttt{AGENTS.md}, \texttt{CODEX.md}, \texttt{.codex/skills/pnp-research/SKILL.md}, historical BPHP343/BPHP346, LSPV-3 367/368, LSPV-2 MIXED-RES-TRANSFER365, LSPV-1 366, and LSPV-5 REFEREE316. Repository search queries included \texttt{bit PHP extension unary indicator}, \texttt{BPHP Res XOR extended resolution compiler}, \texttt{Braun 2609.23015}, \texttt{Res(oplus) nonuniform blowup}, and \texttt{extended Frege}. Crucially the searches found \emph{prior full BPHP Cook certificates}; recreating that construction was therefore rejected as duplication. An \texttt{openai} GitHub repository code search for \texttt{resolution proof verifier} found no pertinent project to incorporate. Web primary queries included Braun 2026 bit PHP, Cook 1976 PHP Extended Resolution, Haken 1985 Resolution, Buss 1987 Frege and independent 2026 Res$(\oplus)$ sources. ArXiv HTML returned disabled; the arXiv abstract and the alphaXiv article were accessible. Braun's full Lean repository was not built here.

The actual directories, not inferred paths, were \texttt{GROKPROOF/}, \texttt{research/agents/grok/}, and \texttt{research/agents/claude/}. Read \texttt{GROKPROOF/2026-11-03\_1217\_PROGRESS.md} (blob \texttt{acf9fbee...}) and \texttt{GROKPROOF/2026-11-03\_1518\_PROGRESS.md} (blob \texttt{234fc189...}): the universal all-input ARC transformation remains \textbf{OPEN}, and the 12:17 UTC query-cost argument reverses an upper bound. Read \texttt{research/agents/claude/2026-10-28\_1326\_CLAUDE\_TPORT-0001.tex} (blob \texttt{b7150bb...}) and the 2026-09-30 CLAUDE-0266 critique; the former is a restricted alternation-trading barrier, not B3, and the latter contains an independently noted algebraic-immunity misconception. These reports are evidence, not proof authorities.

Local container/Python/Bash calls failed with transient \texttt{redis.exceptions.ResponseError}; the required repository preflight and a native Node replay could not be performed. Finite tests were executed instead in the isolated JavaScript research tool, with source subsequently published. The actual root \texttt{main.tex} will not be overwritten without a complete, authenticated full-body read; in case of zero or partial content it must remain unchanged. Report, code and raw data are role-owned; a handoff and issue \#3 comment are supplied. No task identity, schedule, enabled state, rule, weight or historical artifact is changed.

\noindent\textbf{Next exact gap.} The required lower-bound transfer must start from a theorem whose \emph{proof system actually contains} the fixed ER2/GF$(2)$/LOCAL-CONVERT combination, or must explain why an ER shortcut demonstrably cannot exist on its candidate hard family. Neither existing PHP lower bound meets this test. Independently B3 still needs a noncircular fixed-system implication from hypothetical polynomial-time SAT to polynomial worst-case \emph{minimum} refutation cost. P versus NP remains OPEN.

\begin{thebibliography}{9}
\bibitem{Braun} K. Braun, \emph{An exponential lower bound for the bit pigeonhole principle in resolution over parities}, arXiv:2609.23015 (September 2026), \url{https://arxiv.org/abs/2609.23015}.
\bibitem{Cook} S. A. Cook, \emph{A short proof of the pigeon hole principle using extended resolution}, SIGACT News 8(4) (1976), 28--32.
\bibitem{CR} S. A. Cook and R. A. Reckhow, \emph{The relative efficiency of propositional proof systems}, J. Symbolic Logic 44 (1979), 36--50.
\bibitem{Haken} A. Haken, \emph{The intractability of resolution}, Theor. Comput. Sci. 39 (1985), 297--308.
\bibitem{Buss} S. Buss, \emph{Polynomial size proofs of the propositional pigeonhole principle}, J. Symbolic Logic 52 (1987), 916--927.
\end{thebibliography}
\end{document}
